\section{The literal fixed-shift map and its prime-resolved lift}
\label{sec:literal-map}

\subsection{Rows, coefficients, and the terminal-admissible set}

We use the exact left-packet convention of TPC--53
\citep{WangTPC53}.  Fix \(h_0\ne0\), put \(H=\rad|h_0|\), and write a
source row and an opposite row as
\begin{equation}
 \begin{aligned}
 m&=d_m\ell_m,&
 m_\gamma&=d_\gamma\ell_\gamma,\\
 \ell_m,\ell_\gamma&\in[L,2L],&
 d_m,d_\gamma&\in[D,2D],\\
 (d_m,H)&=1,&
 (d_\gamma,H)&=1.
 \end{aligned}
 \label{eq:tpc55-rows}
\end{equation}
with prime \(\ell\)'s and squarefree \(d\)'s.  The inherited separation
makes an integer \(m\) determine \((d_m,\ell_m)\) on the declared cell.
Write
\begin{equation}
 Q=LD,
 \qquad K=DJ,
 \qquad QJ\asymp X,
 \qquad \kappa_0=\frac1{400},
 \qquad \Delta=QX^{-\kappa_0},
 \qquad G=X^{\kappa_0}.
 \label{eq:tpc55-scales}
\end{equation}
Fix one direct row-residue group, one orbit class
\(a\pmod{H_{\rm orb}}\), one opposite-row box \(\cQ\), and one compact
orbit interval \(I_0\Subset(0,\infty)\).  Here
\(H_{\rm orb}=O_{h_0}(1)\) denotes the common
orbit period; it is not being identified with \(|h_0|\).

The weights \(\omega_D\) and \(\psi_L\) range over the fixed bounded
smooth dyadic families inherited from TPC--53: after the displayed
normalization, their supports lie in \([D,2D]\) and \([1,2]\),
respectively, and all normalized derivatives are uniformly bounded.
We use the exact alias \(\zeta_m:=\zeta_m^{(1)}\).  Thus the source
coefficient is
\begin{equation}
 a_m
 =\mu(d_m)(\log\ell_m)\omega_D(d_m)
   \psi_L(\ell_m/L)\zeta_m,
 \qquad |\zeta_m|\ll1.
 \label{eq:tpc55-source-coefficient}
\end{equation}
The opposite row has unary factor
\(\eta_\gamma\mathscr R(q_\gamma)\), where
\begin{equation}
 q_\gamma
 :=\left(\log\frac{\ell_\gamma}{L},
          \log\frac{d_\gamma}{D}\right),
 \qquad |\eta_\gamma|=1,
 \label{eq:tpc55-normalized-opposite-row}
\end{equation}
and \(\mathscr R\) is the fixed smooth function denoted by \(R\) in
TPC--53, relabeled here to avoid collision with the cofactor scale \(R\).
It is supported on the normalized opposite-row
box.  The direct
orbit cell contributes a constant \(\xi_a\ne0\).  The literal binary mask
is
\begin{equation}
 K_{m\gamma}
 =\one_{\ell_m\ne\ell_\gamma}
  \one_{|m-m_\gamma|>\Delta}
  \one_{(d_m,d_\gamma)\le G}.
 \label{eq:tpc55-mask}
\end{equation}
For the declared four-factor profile, set
\begin{equation}
 \cW_{m,\gamma}(j/J)
 =W_1(mj/X)\Psi_1(d_mj/K)
  W_2(m_\gamma j/X)\Psi_2(d_\gamma j/K),
 \label{eq:tpc55-four-factor-profile}
\end{equation}
where \(W_1,W_2,\Psi_1,\Psi_2\in C_c^\infty((0,\infty))\) are the fixed
smooth cutoff functions of the declared packet.

Let \(\cI_X\) be the retained set of pairs \((m,r)\), let \(\cO_X\) be
the complete opposite-row cell, and put
\begin{equation}
 \cJ_{X,a}
 :=\{j:j/J\in I_0,\ j\equiv a\pmod{H_{\rm orb}}\}.
 \label{eq:tpc55-retained-orbit-set}
\end{equation}
The full pre-terminal atomic universe is
\begin{equation}
 \cU^{\rm pre}
 :=\{(m,r,j,\gamma):(m,r)\in\cI_X,\
       j\in\cJ_{X,a},\ \gamma\in\cO_X\}.
 \label{eq:tpc55-preterminal-universe}
\end{equation}
Thus the scalar coordinate of the lifted physical vector is
\begin{equation}
 g_u:=[G_X^L(m,r,j)]_\gamma
 =a_mK_{m\gamma}\eta_\gamma\mathscr R(q_\gamma)\xi_a
   \cW_{m,\gamma}(j/J),
 \qquad u=(m,r,j,\gamma).
 \label{eq:tpc55-literal-atom}
\end{equation}
The value is independent of \(r\) except for repetition over the retained
cofactor support.

Let \(\Phi_0\) be the inherited terminal cutoff and define
\begin{equation}
 \Lambda(n):=(\log n)\Phi_0(n/X).
 \label{eq:tpc55-terminal-Lambda}
\end{equation}
This notation is a shaped prime amplitude, not the von Mangoldt function.
For \(u=(m,r,j,\gamma)\in\cU^{\rm pre}\), put
\begin{equation}
 p_u:=\frac{mj+h_0}{r}
 \label{eq:tpc55-terminal-prime}
\end{equation}
whenever the quotient is integral.  The terminal-admissible universe is
the inherited coprime squarefree equality support
\begin{equation}
 \begin{split}
 \cU_{h_0}^{\rm term}:=\{u\in\cU^{\rm pre}:{}&
 p_u\text{ is a prime in the inherited terminal support},\\
 &p_ur-mj=h_0,\quad \Lambda(p_ur)\ne0,\quad
 d_m,r,d_mr\text{ are squarefree},\quad(d_m,r)=1\}.
 \end{split}
 \label{eq:tpc55-terminal-universe}
\end{equation}
Define separately the actually nonzero physical support
\begin{equation}
 \cU_{h_0}^{\rm phys}
 :=\{u\in\cU_{h_0}^{\rm term}:g_u\ne0\}.
 \label{eq:tpc55-physical-active-set}
\end{equation}
This distinction prevents pre-terminal rows from being silently treated as
terminal-admissible atoms.

Thus, on every terminal-admissible atom, the inherited TPC--45/48 support
conditions \citep{WangTPC45,WangTPC48} give
\begin{equation}
 d_m, r, s=d_mr\ \text{squarefree},
 \qquad (d_m,r)=1,
 \qquad \mu(d_m)\mu(r)=\mu(s).
 \label{eq:tpc55-squarefree-support}
\end{equation}
The complete physical output key and its prime refinement are
\begin{equation}
 q(u)=(\gamma,d_mr,j),
 \qquad
 \widetilde q(u)=(\gamma,d_mr,j,p_u).
 \label{eq:tpc55-two-output-maps}
\end{equation}

\subsection{Three exact operators}

Use counting measure on every atomic coordinate.  For an arbitrary array
\(g=(g_u)_{u\in\cU^{\rm pre}}\), the terminal diagonal is an operator on
the full pre-terminal universe, defined by
\begin{equation}
 (D_{h_0}g)_u
 :=
 \begin{cases}
  \mu(r)\Lambda(p_ur)g_u,
   &u=(m,r,j,\gamma)\in\cU_{h_0}^{\rm term},\\
  0,&u\in\cU^{\rm pre}\setminus\cU_{h_0}^{\rm term}.
 \end{cases}
 \label{eq:tpc55-terminal-diagonal}
\end{equation}
Let \(\widetilde\cY_{h_0}=\widetilde q(\cU_{h_0}^{\rm term})\).  The
prime-resolved lift is
\begin{equation}
 (\cR_{h_0}g)_{\widetilde y}
 :=\sum_{\substack{u\in\cU_{h_0}^{\rm term}\\
                    \widetilde q(u)=\widetilde y}}
       (D_{h_0}g)_u,
 \qquad \widetilde y\in\widetilde\cY_{h_0}.
 \label{eq:tpc55-prime-resolved-lift}
\end{equation}
Finally, let \(\Sigma_p\) forget the prime coordinate:
\begin{equation}
 (\Sigma_p x)_{\gamma,s,j}
 :=\sum_{p:(\gamma,s,j,p)\in\widetilde\cY_{h_0}}
 x_{\gamma,s,j,p}.
 \label{eq:tpc55-forget-prime}
\end{equation}
The exact terminal map of TPC--53 is
 \begin{equation}
  (\cT_{h_0}g)_{\gamma,s,j}
  =\sum_{\substack{u=(m,r,j,\gamma)\in\cU_{h_0}^{\rm term}\\
                    s=d_mr}}
    \mu(r)\Lambda(p_ur)g_{m,r,j,\gamma}.
  \label{eq:tpc55-literal-terminal-map}
 \end{equation}
 \begin{equation}
  \cT_{h_0}=\Sigma_p\cR_{h_0}.
  \label{eq:tpc55-terminal-factorization}
 \end{equation}
Every sum is restricted to the declared finite pre-terminal, terminal,
and output supports.  No projective layer or frozen intercept is present
in these definitions.

\subsection{Prime-resolved injectivity}

The factorization becomes useful because \(\cR_{h_0}\) has no collision
after restriction to the terminal-admissible coordinate space.

\begin{lemma}[Fixed-coordinate fused-label injectivity]
\label{lem:tpc55-fused-injectivity}
Assume \(D=o(J)\), \(h_0\ne0\), and the row uniqueness in
\eqref{eq:tpc55-rows}.  For all sufficiently large \(X\), if two rows in
the same source system satisfy
\begin{equation}
 d_1(d_1\ell_1j+h_0)=d_2(d_2\ell_2j+h_0)
 \label{eq:tpc55-fused-label-equality}
\end{equation}
at the same retained \(j\asymp J\), then
\((d_1,\ell_1)=(d_2,\ell_2)\).
\end{lemma}

\begin{proof}
Rearranging gives
\begin{equation}
 j(\ell_1d_1^2-\ell_2d_2^2)=h_0(d_2-d_1).
 \label{eq:tpc55-fused-injectivity-equation}
\end{equation}
If the integer in parentheses on the left is nonzero, the left side has
absolute value at least a fixed multiple of \(J\), while the right side is
\(O_{h_0}(D)=o(J)\), a contradiction.  Hence both sides vanish.  Since
\(h_0\ne0\), first \(d_1=d_2\), and then \(\ell_1=\ell_2\).
\end{proof}

\begin{theorem}[Prime-resolved desingularization]
\label{thm:tpc55-prime-resolved-injection}
Assume \(h_0\ne0\), \(D=o(J)\), uniqueness of the row decomposition in
\eqref{eq:tpc55-rows}, and sufficiently large \(X\).  On one declared
source system, the map
\(\widetilde q:\cU_{h_0}^{\rm term}\to\widetilde\cY_{h_0}\) is
injective.  Consequently the restriction
\(\cR_{h_0}|_{\ell^2(\cU_{h_0}^{\rm term})}\) is a weighted coordinate
embedding.  After the separate restriction from \(\cU^{\rm pre}\) to
\(\cU_{h_0}^{\rm term}\), all many-to-one behavior of \(\cT_{h_0}\)
occurs in the forget-prime map \(\Sigma_p\).
\end{theorem}

\begin{proof}
Suppose \(\widetilde q(u_1)=\widetilde q(u_2)\).  Then the two atoms have
the same \(s,j,p,\gamma\).  Their fused labels satisfy
\[
 d_i(d_i\ell_i j+h_0)
 =d_i(p r_i)=p(d_ir_i)=ps.
\]
Thus \eqref{eq:tpc55-fused-label-equality} holds, and
\cref{lem:tpc55-fused-injectivity} gives the same source row.  Then
\(r_i=s/d_i\) is also the same, so \(u_1=u_2\).  Formula
\eqref{eq:tpc55-terminal-factorization} shows that only \(\Sigma_p\) can
identify distinct inputs.
\end{proof}

\begin{corollary}[Terminal primes parameterize a full collision fiber]
\label{cor:tpc55-fiber-prime-bijection}
For \(y=(\gamma,s,j)\), put
\begin{equation}
 \begin{aligned}
 \cF_y^{\rm term}&:=q^{-1}(y)\cap\cU_{h_0}^{\rm term},&
 \cF_y^{\rm phys}&:=q^{-1}(y)\cap\cU_{h_0}^{\rm phys},\\
 \cP_y^{\rm term}&:=\{p_u:u\in\cF_y^{\rm term}\},&
 \cP_y^{\rm phys}&:=\{p_u:u\in\cF_y^{\rm phys}\}.
 \end{aligned}
 \label{eq:tpc55-fiber-and-primes}
\end{equation}
For compatibility below, write
\(\cF_y:=\cF_y^{\rm term}\) and
\(\cP_y:=\cP_y^{\rm term}\).  Then \(u\mapsto p_u\) is a bijection from
\(\cF_y^{\rm term}\) to \(\cP_y^{\rm term}\), and its restriction is a
bijection from \(\cF_y^{\rm phys}\) to \(\cP_y^{\rm phys}\).  In
particular, two distinct terminal-admissible atoms in one full output
fiber carry distinct terminal primes.
\end{corollary}

\subsection{The three energies}

The preceding factorization keeps three norms separate.  For an arbitrary
lifted array \(g\), define
\begin{align}
 \cE_{\rm pre}(g)
 &:=\sum_{u\in\cU^{\rm pre}}|g_u|^2,
 \label{eq:tpc55-preterminal-energy}\\
 \cD_{h_0}(g)
 &:=\|D_{h_0}g\|_2^2
 =\sum_{u\in\cU_{h_0}^{\rm term}}
   |\mu(r)|^2|\Lambda(p_ur)|^2|g_u|^2\notag\\
 &=\sum_{u\in\cU_{h_0}^{\rm term}}
   |\Lambda(p_ur)|^2|g_u|^2,
 \label{eq:tpc55-terminal-diagonal-energy}\\
 \cG_{h_0}(g)
 &:=\|\cT_{h_0}g\|_2^2
 =\sum_y\left|\sum_{u\in\cF_y^{\rm term}}
   \mu(r_u)\Lambda(p_ur_u)g_u\right|^2.
 \label{eq:tpc55-grouped-energy}
\end{align}
The last equality in \(\cD_{h_0}\) uses the inherited squarefree
condition \(r\) on \(\cU_{h_0}^{\rm term}\), so
\(|\mu(r)|^2=1\); it would be false on an unrestricted cofactor universe.
The prime-resolved injection gives the exact identity
\begin{equation}
 \|\cR_{h_0}g\|_2^2=\cD_{h_0}(g).
 \label{eq:tpc55-resolved-isometry}
\end{equation}
It does not give \(\cD_{h_0}\asymp\cE_{\rm pre}\): that comparison also
requires terminal activation and two-sided control of \(\Lambda\) on the
selected support.  Nor does it give a lower bound for
\(\cG_{h_0}/\cD_{h_0}\), because \(\Sigma_p\) creates the cross terms.

For the physical vector \(G_X^L\), the pre-terminal energy is the raw
TPC--53 lifted profile sum.  It should not be confused with the normalized
\(1/(E_XJ)\) model statistic used at an earlier intermediate stage in
TPC--52 \citep{WangTPC52}.  Equations
\eqref{eq:tpc55-preterminal-energy}--\eqref{eq:tpc55-grouped-energy} use
literal counting-measure atoms throughout.
